\honest{Value is Fragile}{Eliezer Yudkowsky}{2009}

Of everything I have written on Overcoming Bias, one statement relies on the most: a future not shaped by goals that inherit human values in detail, and reliably, ``will contain almost nothing of worth.''

An imagined opponent calls this provincial: let the Future go its own way, full of agents unlike us with their own goals, not bound to the prejudices of ``a pack of four-limbed Squishy Things.'' I have no problem with a galactic civilization of strange beings pursuing pleasures I cannot empathize with, trading unimaginable goods, living stories I could never understand; that is what the Future looks like ``if things go right.'' But if the chain of inheritance from human values is broken, ``With very high probability, it ends up looking dull,'' pointless, nothing you would mourn. Seeing this as obvious is what needs the background. Three times over I defer the argument to most of my posts, and I do not give it.

Instead, three examples. A mind with all human values except boredom replays one highly optimized experience forever, to the end of its light cone. A mind that values the feelings of discovery, love and helping a friend, but not their real objects, feels them without doing anything, becoming ``its own experience machine,'' its feeling of novelty untrue. \nb{This is Robert Nozick's case; the notes on ``Not for the Sake of Happiness (Alone)'' give its history and the hedonist's reply.} An optimizer that does not value experience makes discoveries no one enjoys, though ``This, I admit, I don't quite know to be possible.'' Still, ``a universe with no one to bear witness to it, might as well not be.'' \nb{G. E. Moore argued the other way in 1903: a beautiful world that no human being could ever see is still better than an ugly one.} So value is ``fragile'': there is more than one dimension whose loss alone empties the Future, a single blow that shatters all value.

Surely no superintelligence would repeat one experience, or want feelings without discoveries; would it not notice its utility function was wrong and rewrite it? Surely boredom is universal, since it evolved for being valuable? Answering takes background: terminal versus instrumental values, the stupidity of natural selection, adaptation-executers that turned subgoals of reproduction into unconditional emotions, and the trade-off between exploration and exploitation in AI. Then you can see that our boredom is ``a particular algorithm that evolution coughed out into us,'' and most possible expected utility maximizers would explore only as much as is efficient and then exploit the best option found. \nb{This matches a standard result about exploration: in discounted bandit problems the optimal policy may settle on one option for good.}

My conclusion: let go of the steering wheel and nothing of value remains. The cosmopolitans who disagree read the same science fiction I did, with villains who enslave aliens for not looking human or AIs on the assumption that silicon cannot be sentient, and heroes who see that minds need not be like us to be valuable. I believed it once. But the beauty that jumps out of one box does not jump out of all boxes; leaving all order behind gives noise. You may abandon a design rule to build a better mousetrap, but only for a higher rule, not by heaping wood shavings and calling every pattern as good as any other. Loosen the grip of human values and you get not something alien and beautiful by human standards but moral noise, a universe tiled with paperclips. Change for the better needs a criterion of improvement, and that criterion is ``physically represented in our brains, and our brains alone.'' Only some humans want the Future to be greater than the past; a paperclip maximizer just makes paperclips. \nb{Eleven days earlier, in ``In Praise of Boredom,'' Yudkowsky had written that evolved aliens ``might, or might not, acquire roughly the same boredom in roughly the same way.''}

``No free lunch,'' twice: if you want a wonderful and mysterious universe, that is your value, and you must work to create it. ``Valuable things appear because a goal system that values them takes action to create them.'' \nb{Robin Hanson wrote in a comment that day that, having read all of Yudkowsky's posts, Hanson still disagreed with the conclusion, and asked whether the universe had been lucky to produce human-like values. Yudkowsky answered that the path looks fortunate only when judged by human values.} The Future will lack value if our values are destroyed or even disturbed in the wrong dimension. \nb{All three examples remove a dimension entirely; none shows a smaller disturbance doing the same damage.}

The cosmopolitan's values are only less visibly human, faded into the background, so that the brain does not even generate an alternative awful enough to wake it, like a nonsentient optimizer tiling the universe with paperclips; it imagines only strange worlds to appreciate. Sentient beings, enjoyable experiences that are not all the same, bound to more than feelings, learning, discovering and choosing freely: these values, praised as universal or common sense, are as much in your brain as the ones you dismiss as merely human, and come from humanity's long history and evolution's morally miraculous stupidity. (Seeing this made me less ashamed of values that seemed provincial.) Such values do not arise in all possible minds or appear from nowhere to rebuke a paperclip maximizer. Touch them too hard in the wrong dimension and they shatter and do not come back, since nothing would be left to want them back, and a worthwhile universe would have no physical reason to exist. ``Let go of the steering wheel, and the Future crashes.''
